\documentclass{article}

% Recommended, but optional, packages for figures and better typesetting:
\usepackage{microtype}
\usepackage{graphicx}
\usepackage{subcaption}
\usepackage{booktabs} % for professional tables
\usepackage{hyperref}
\newcommand{\theHalgorithm}{\arabic{algorithm}}
\usepackage{icml2026}
\usepackage{algorithm}
\usepackage{algorithmic}

\usepackage{amsmath}
\usepackage{amssymb}
\usepackage{mathtools}
\usepackage{amsthm}
\usepackage[capitalize,noabbrev]{cleveref}

\theoremstyle{plain}
\newtheorem{theorem}{Theorem}[section]
\newtheorem{proposition}[theorem]{Proposition}
\newtheorem{lemma}[theorem]{Lemma}
\newtheorem{corollary}[theorem]{Corollary}
\theoremstyle{definition}
\newtheorem{definition}[theorem]{Definition}
\newtheorem{assumption}[theorem]{Assumption}
\theoremstyle{remark}
\newtheorem{remark}[theorem]{Remark}

\icmltitlerunning{Self-Refining Calibration-Free Prototype Alignment for Test-Time Model Merging}

\begin{document}

\twocolumn[
  \icmltitle{Self-Refining Calibration-Free Prototype Alignment for \\ Robust Test-Time Model Merging}

  \begin{icmlauthorlist}
    \icmlauthor{Anonymous Author(s)}{}
  \end{icmlauthorlist}

  \icmlaffiliation{dept}{Department of Computer Science, Anonymous University, Anonymous Country}
  \icmlcorrespondingauthor{Anonymous}{anonymous@anonymous.edu}

  \icmlkeywords{Model Merging, Test-Time Adaptation, Self-Supervised Learning, Prototype Alignment}

  \vskip 0.3in
]

\printAffiliationsAndNotice{}

\begin{abstract}
Test-Time Model Merging (TTMM) is an emerging paradigm that dynamically fuses multiple task-specific expert models to adapt to a non-stationary stream of unlabeled test data. However, existing methods either rely on offline calibration sets to pre-compute clean prototypes (which suffer from severe semantic mismatch under environmental corruptions) or suffer from softmax saturation and representation collapse during rapid task transitions. To address these limitations, we propose \textbf{Self-Refining Calibration-Free Prototype Alignment (SR-CPA)}, a fully calibration-free and training-data-free TTMM framework. SR-CPA initializes class prototypes on-the-fly directly from the incoming test stream and continuously refines them using an exponential moving average (EMA) of high-confidence representations. This self-refining mechanism naturally compensates for domain drift and aligns prototypes directly with the corrupted feature space. Under non-stationary test streams with severe corruptions, SR-CPA demonstrates superior adaptability, achieving an average accuracy of \textbf{13.29\%} (and up to \textbf{13.62\%} with uncertainty-aware adaptive temperatures and Fisher-weighted learning rates) on rapid task-alternating streams, outperforming existing calibration-based and gradient-based adaptation baselines by a large margin while matching or exceeding performance on sequential streams.
\end{abstract}

\section{Introduction}
\label{sec:intro}

Deep neural networks deployed in real-world environments frequently encounter non-stationary data streams characterized by sudden shifts in both task semantics and environmental conditions (e.g., weather corruptions, noise, blur, and lighting changes). Fully adjusting a massive multi-task network or keeping multiple specialized expert models active in memory at test-time is often computationally prohibitive. Consequently, \textbf{Test-Time Model Merging (TTMM)} \cite{cpamerge2025, lfwa2025} has emerged as a promising, lightweight alternative. In TTMM, specialized expert models (e.g., pre-trained on different tasks) are dynamically merged into a single student model at test-time using layer-wise merging coefficients $\boldsymbol{\Lambda}$, without storing or backpropagating through the original teacher networks.

Despite its potential, deploying TTMM in wild, unconstrained test streams remains highly challenging. First, standard Test-Time Adaptation (TTA) approaches \cite{wang2020tent, niu2023towards} optimize merging coefficients via unsupervised entropy minimization, which is highly prone to softmax saturation and representation collapse. Once coefficients saturate to a single expert model, the student network loses the capability to transition or adapt to subsequent tasks in the stream. Second, state-of-the-art TTMM methods like \textbf{CPA-Merge} \cite{cpamerge2025} attempt to prevent this collapse by pre-computing static class prototypes on offline, clean calibration datasets to guide dynamic routing. However, under non-stationary streams with environmental corruptions (e.g., severe Gaussian noise, motion blur, and contrast shifts), these static, clean prototypes exhibit a severe \textbf{semantic mismatch} with the corrupted representation space. This mismatch leads to incorrect routing decisions and catastrophic failure under rapid task switches.

To resolve these challenges, we introduce \textbf{Self-Refining Calibration-Free Prototype Alignment (SR-CPA)}, a novel, fully calibration-free test-time model merging framework. SR-CPA eliminates the dependency on offline calibration data or clean training sets entirely by initializing and refining class prototypes *on-the-fly* directly from the corrupted test stream. 

Our framework operates by first leveraging the initial uniform merged model as a stable feature anchor to extract representations from the early batches of the stream. It filters for high-confidence predictions ($\mathcal{H}_{thresh} > 0.90$) to seed initial class prototypes. As adaptation progresses under non-stationary environmental corruptions, SR-CPA continuously refines and consolidates these prototypes via an exponential moving average (EMA) of newly observed high-confidence embeddings. By updating the prototypes directly within the corrupted representation space, SR-CPA naturally compensates for domain drift and ensures that the prototypes remain highly aligned with the incoming test features.

Through extensive simulation experiments, we evaluate SR-CPA against state-of-the-art TTMM baselines across a variety of streams and corruptions (clean, noise, blur, contrast). Our results show that:
\begin{itemize}
    \item On sequential streams under corruptions, SR-CPA performs comparably to calibration-based prototype alignment, while running in a completely training-data-free manner.
    \item On rapid task-alternating streams, our proposed methods achieve a remarkable average accuracy of \textbf{13.29\%} (and up to \textbf{13.62\%} with uncertainty-aware adaptive temperatures and Fisher-weighted learning rates), which actually exceeds the uniform STATIC baseline (\textbf{13.48\%}), while CPA-Merge collapses completely to \textbf{7.51\%}. On the clean alternating stream, SR-CPA reaches \textbf{16.62\%} (close to STATIC's 16.69\%) while CPA-Merge collapses to \textbf{3.42\%}. Additionally, on sequential streams, our adaptive temperature variant (\textbf{AT-SR-CPA}) achieves the highest overall average sequential accuracy of \textbf{12.94\%}, outperforming the standard unconstrained TTA baseline (\textbf{12.86\%}).
    \item SR-CPA maintains high stability and is highly robust to hyperparameter selection, as shown by comprehensive sensitivity sweeps over the softmax routing temperature $\tau$ and the prototype EMA update rate $\gamma$.
\end{itemize}

\section{Related Work}
\label{sec:related}

\textbf{Model Merging:} Weight averaging and model merging have gained traction as a way to combine multiple specialized neural networks into a single multi-task network without additional training. Early works merged models via simple weight averaging \cite{wortsman2022model} or permutation alignment \cite{ainsworth2022git}. Subsequent methods incorporated structural priors such as Fisher Information \cite{matena2022merging}, class-specific gradient/activation covariance \cite{jin2022regmean, stoica2023zipit}, and interference-reduction techniques like Ties-Merging \cite{yadav2023resolving} and DARE \cite{yu2023language}. While these methods assume a static multi-task setting, our work focuses on the more challenging non-stationary test-time setting.

\textbf{Test-Time Adaptation (TTA):} TTA aims to adapt a pre-trained model to out-of-distribution (OOD) test data without accessing the source training data. TENT \cite{wang2020tent} minimizes the entropy of model predictions at test-time. Continual TTA methods \cite{wang2022continual, gong2022note, niu2023towards} introduce regularization and sample-selection strategies to prevent catastrophic forgetting and error accumulation over long streams. However, standard TTA methods generally optimize the entire network weight space, which incurs high computational overhead and risks representation collapse under extreme noise. Our framework, in contrast, optimizes only a few layer-wise merging coefficients, which is highly parameter-efficient and stable.

\textbf{Test-Time Model Merging (TTMM):} Combining model merging and test-time adaptation, TTMM seeks to adaptively merge multiple task-specific expert models at test-time based on the incoming stream. Recent works have explored various facets of TTMM: for instance, \cite{bertolissi2025local} show that TTMM is a highly efficient approximation of test-time training (TTT) that scales to numerous experts, and \cite{qiu2025mingle} introduce MINGLE to handle test-time continual model merging via null-space gating. 
Other methods like LFWA \cite{lfwa2025} use diagonal Fisher Information as a static sensitivity prior to scale layer-wise learning rates dynamically, which stabilizes adaptation but requires expert training data. PC-Merge \cite{pcmerge2025} implements gradient surgery and hard parameter resets to handle abrupt task transitions, which requires expensive pairwise gradient projections. CPA-Merge \cite{cpamerge2025} introduces prototype-driven dynamic routing (PD-Routing) and contrastive alignment using clean, pre-computed prototypes. Unlike these methods, SR-CPA is completely calibration-free and dynamically refines its prototypes to track corrupted and non-stationary domains without any offline training or calibration data, or complex low-rank gating parameters.

\section{Proposed Method: SR-CPA}
\label{sec:method}

We consider a multi-task test-time model merging setting with $K$ specialized expert models, where each expert $k \in \{0, \ldots, K-1\}$ has a set of pre-trained parameters $\boldsymbol{W}^{(k)}$. At test-time, the student model parameters are represented as a layer-wise convex combination of a base model $\boldsymbol{W}_{base}$ and the task-specific expert deviations (the "task vectors") $\boldsymbol{\tau}_l^{(k)} = \boldsymbol{W}_l^{(k)} - \boldsymbol{W}_{base, l}$:
\begin{equation}
    \boldsymbol{W}_l(\boldsymbol{\Lambda}_l) = \boldsymbol{W}_{base, l} + \sum_{k=0}^{K-1} \Lambda_{l,k} \boldsymbol{\tau}_l^{(k)}
\end{equation}
where $l \in \{0, \dots, L-1\}$ indexes the layers of the model, and $\boldsymbol{\Lambda}_l = [\Lambda_{l,0}, \dots, \Lambda_{l,K-1}]^T$ represents the merging coefficients for layer $l$, constrained to lie on the probability simplex $\Delta^{K-1}$.

\subsection{Calibration-Free Prototype Initialization}
\label{subsec:init}

Unlike existing prototype-based merging methods that require pre-computing prototypes offline on clean training data, SR-CPA initializes class prototypes on-the-fly directly from the unlabeled test stream. 

We initialize a set of self-refined prototypes $\boldsymbol{P} = \{\boldsymbol{p}_{k,c}\}$, where $k \in \{0, \ldots, K-1\}$ represents the task, $c \in \{0, \ldots, C-1\}$ indexes the classes, and $\boldsymbol{p}_{k,c} \in \mathbb{R}^D$ is the prototype vector. At the beginning of the stream, all prototypes are initialized to zero.

For each incoming test batch $B_t = \{\boldsymbol{x}_i\}$, we first compute a forward anchor pass through the static uniform merged model ($\boldsymbol{\Lambda}_l = [1/K, \ldots, 1/K]^T$) to extract intermediate representation features $\boldsymbol{z}_i^{anchor}$ and predict class probabilities $\boldsymbol{\hat{y}}_i = \text{Softmax}(\boldsymbol{z}_i^{anchor} \boldsymbol{H}_k^T)$, where $\boldsymbol{H}_k$ is the classification head for task $k$. We filter for high-confidence predictions where the maximum class probability exceeds a threshold $\mathcal{H}_{thresh} = 0.90$.

For each high-confidence sample $i$, we extract its representation embedding and assign it to the predicted class $c = \text{argmax}(\boldsymbol{\hat{y}}_i)$ for the active task $k$. If the prototype $\boldsymbol{p}_{k,c}$ is uninitialized (i.e., all zeros), we directly seed it with the normalized feature representation:
\begin{equation}
    \boldsymbol{p}_{k,c} = \text{Normalize}\left(\boldsymbol{z}_i^{anchor}\right)
\end{equation}
where $\text{Normalize}(\boldsymbol{u}) = \boldsymbol{u} / \|\boldsymbol{u}\|_2$. This simple initialization enables the student model to construct class centroids directly within the active domain representation space.

\subsection{Dynamic Self-Refining Prototype Update}
\label{subsec:refine}

Under non-stationary test streams, environmental corruptions (like blur or contrast shifts) cause the feature representation space of incoming samples to drift continuously. Static pre-computed prototypes suffer from a severe domain mismatch under these shifts. To compensate for this drift, SR-CPA dynamically refines and consolidates the prototypes during test-time adaptation.

For each high-confidence sample $i$ in batch $B_t$ with predicted class $c = \text{argmax}(\boldsymbol{\hat{y}}_i)$, we update the existing prototype $\boldsymbol{p}_{k,c}$ via an exponential moving average (EMA) with an update rate $\gamma_i \in (0, 1)$:
\begin{align}
    \boldsymbol{p}_{k,c}^{(t)} &= \text{Normalize}\left( (1 - \gamma_i) \boldsymbol{p}_{k,c}^{(t-1)} + \gamma_i \boldsymbol{z}_{i}^{(t)} \right)
\end{align}
where $\boldsymbol{z}_{i}^{(t)}$ is the representation embedding of the sample extracted under the *currently active, adapted* merging coefficients. By default, $\gamma_i = \gamma$ is a constant update rate. 

To further enhance robustness and prioritize highly reliable representations, we also propose a novel variant: \textbf{Confidence-Weighted Self-Refinement (CW-SR-CPA)}, where the prototype update rate is dynamically scaled by the prediction confidence of each individual sample:
\begin{equation}
    \gamma_i = \gamma \times \max \hat{\boldsymbol{y}}_i
\end{equation}
This formulation ensures that highly clear and confident predictions contribute more to the prototype definitions, whereas less confident predictions (which are closer to the threshold and more likely to represent noise or boundary classes) have a mitigated impact, preventing prototype pollution and error propagation.

By using adapted representations to refine the prototypes, SR-CPA ensures that the prototypes are actively aligned with the adapted representation space. This self-refining loop provides powerful drift-compensation and stabilizes the semantic definitions of class prototypes over long and noisy streams.

\subsection{Theoretical Convergence and Stability Analysis}
\label{subsec:theory}

To mathematically justify the effectiveness and stability of the Self-Refining Prototype Alignment framework, we analyze the convergence behavior of our prototype update rule. We assume that the feature representation $\boldsymbol{z}_t^{(c)}$ of a sample belonging to class $c$ of task $k$ at step $t$ in the corrupted domain is given by:
\begin{equation}
    \boldsymbol{z}_t^{(c)} = \boldsymbol{\mu}_{k,c} + \boldsymbol{\epsilon}_t
\end{equation}
where $\boldsymbol{\mu}_{k,c}$ is the true domain-specific centroid of class $c$, and $\boldsymbol{\epsilon}_t$ is a zero-mean i.i.d.\ noise vector representing representation variance under environmental corruptions, satisfying $\mathbb{E}[\boldsymbol{\epsilon}_t] = \boldsymbol{0}$ and $\mathbb{E}[\|\boldsymbol{\epsilon}_t\|_2^2] = \sigma^2$.

Due to classification errors during adaptation, the class assignment for a high-confidence sample is correct with probability $p$ (which we denote as prototype purity) and incorrect with probability $1-p$. Let incorrect samples be drawn from other classes with average representation centroid $\boldsymbol{\mu}_{\text{other}} \neq \boldsymbol{\mu}_{k,c}$. Under these assumptions, we characterize the bias and variance of the unnormalized self-refined prototypes.

\begin{theorem}[Bias-Variance Properties of Self-Refining Prototypes]
Let $\tilde{\boldsymbol{p}}_{k,c}^{(t)}$ be the unnormalized self-refined prototype for class $c$ of task $k$ at step $t$, updated via the linear EMA operator:
\begin{equation}
    \tilde{\boldsymbol{p}}_{k,c}^{(t)} = (1 - \gamma) \tilde{\boldsymbol{p}}_{k,c}^{(t-1)} + \gamma \boldsymbol{z}_t
\end{equation}
The asymptotic expectation and variance of $\tilde{\boldsymbol{p}}_{k,c}^{(t)}$ satisfy:
\begin{align}
    \lim_{t \to \infty} \mathbb{E}[\tilde{\boldsymbol{p}}_{k,c}^{(t)}] &= p \boldsymbol{\mu}_{k,c} + (1-p)\boldsymbol{\mu}_{\text{other}} \label{eq:theory_bias} \\
    \lim_{t \to \infty} \text{Var}(\tilde{\boldsymbol{p}}_{k,c}^{(t)}) &= \frac{\gamma}{2 - \gamma} \text{Var}(\boldsymbol{z}_t) \label{eq:theory_var}
\end{align}
\end{theorem}

\begin{proof}
See Appendix~\ref{sec:proofs} for the full mathematical proof.
\end{proof}

\textbf{Key Theoretical Insights:}
Equation~\eqref{eq:theory_bias} proves that the prototype centroid converges exactly to the true corrupted class center $\boldsymbol{\mu}_{k,c}$ when purity $p \to 1$. Our empirical seeding threshold sweep in Section~\ref{subsec:sensitivity} validates this: setting a conservative $\mathcal{H}_{thresh} \ge 0.85$ keeps purity extremely high ($p \approx 1$), making the prototypes highly representative. Lower thresholds (e.g., $\mathcal{H}_{thresh}=0.70$) introduce severe label noise ($p \ll 1$), which biases the prototype toward incorrect classes and triggers catastrophic collapse.

Equation~\eqref{eq:theory_var} demonstrates the powerful variance-damping effect of our EMA refinement. For a default update rate of $\gamma = 0.1$, the steady-state variance of the refined prototype is $\frac{0.1}{1.9} \approx 5.26\%$ of the raw sample variance. This significant variance reduction explains why SR-CPA's prototypes remain highly stable and completely immune to transient representation noise, preventing representation collapse even under severe corruptions.

\subsection{Softmax Routing Prior (PD-Routing)}
\label{subsec:routing}

To adapt the layer-wise merging coefficients $\boldsymbol{\Lambda}_l$ dynamically for each incoming batch, we leverage the self-refined prototypes to compute a task affinity score. 

For an incoming batch, we first compute the normalized anchor features $\boldsymbol{z}_i^{anchor}$. We then measure the similarity of these features to the prototypes of each task $k$ as the average of the maximum cosine similarities over the batch:
\begin{equation}
    s_k = \frac{1}{|B|} \sum_{i \in B} \max_{c \in \{0,\dots,C-1\}} \cos\left(\boldsymbol{z}_i^{anchor}, \boldsymbol{p}_{k,c}\right)
\end{equation}
If the prototypes for task $k$ are not yet initialized, we assign a neutral affinity score of $0$. We then apply a sharp softmax routing prior with a temperature $\tau_{temp}$ to obtain the routing coefficients:
\begin{equation}
    \text{Prior}_k = \frac{\exp(s_k / \tau_{temp})}{\sum_{j=0}^{K-1} \exp(s_j / \tau_{temp})}
\end{equation}
This prior represents our current confidence in the active task of the batch.

To prevent incorrect, hard routing decisions under high noise or extreme OOD corruptions where prototypes might not be fully formed, we also propose a novel, highly robust variant: \textbf{Uncertainty-Aware Adaptive Temperature (AT-SR-CPA)}. Under this variant, the routing temperature $\tau_{temp}$ is adapted dynamically based on the entropy (uncertainty) of the task affinities. First, we compute a baseline task probability distribution $q$ using a moderate baseline temperature $\tau_0 = 0.05$:
\begin{equation}
    q_k = \frac{\exp(s_k / \tau_0)}{\sum_{j=0}^{K-1} \exp(s_j / \tau_0)}
\end{equation}
We then measure the normalized uncertainty $U \in [0, 1]$ as the ratio of the Shannon entropy of $q$ to the maximum possible task entropy:
\begin{equation}
    U = \frac{-\sum_{k=0}^{K-1} q_k \log q_k}{\log K}
\end{equation}
Using this normalized uncertainty, we dynamically modulate the routing temperature linearly between a minimum temperature $\tau_{min} = 0.01$ and a maximum temperature $\tau_{max} = 0.50$:
\begin{equation}
    \tau_{temp}^{dynamic} = \tau_{min} + (\tau_{max} - \tau_{min}) \times U
\end{equation}
Under high task certainty ($U \approx 0$), the temperature remains extremely sharp ($\approx 0.01$), routing the model decisively to the active task and pruning interference. Conversely, under high domain uncertainty ($U \approx 1$), the temperature becomes highly diffuse ($\approx 0.50$), generating a near-uniform routing prior. This automatically cushions the model against erroneous task-tracking decisions, falling back gracefully to stable uniform routing.

We reset the merging coefficients $\boldsymbol{\Lambda}_l$ across all layers directly to this prior prior to backpropagation:
\begin{equation}
    \boldsymbol{\Lambda}_l \leftarrow \textbf{Prior} \quad \forall l \in \{0, \dots, L-1\}
\end{equation}
By resetting $\boldsymbol{\Lambda}_l$ to this dynamic routing prior at each step, we completely bypass the "softmax saturation" trap of TTA, allowing the coefficients to remain agile and adapt instantaneously to abrupt task switches.

\subsection{Contrastive Prototype Alignment Loss}
\label{subsec:loss}

To further optimize the merging coefficients and align the representations, we apply a confidence-masked contrastive alignment loss (InfoNCE) using the self-refined prototypes.

For high-confidence samples where the maximum predicted probability exceeds $\mathcal{H}_{contra} = 0.85$, we define the contrastive similarity matrix $\boldsymbol{M}$ between the normalized adapted representation $\boldsymbol{z}_i$ and the normalized self-refined prototypes of the active task $\boldsymbol{p}_{active, c}$:
\begin{equation}
    M_{i,c} = \frac{\cos\left(\boldsymbol{z}_i, \boldsymbol{p}_{active, c}\right)}{\kappa}
\end{equation}
where $\kappa = 0.1$ is the contrastive temperature. The InfoNCE loss over high-confidence samples is given by:
\begin{equation}
    \mathcal{L}_{contra} = -\frac{1}{|B_{conf}|} \sum_{i \in B_{conf}} \log \frac{\exp(M_{i, \hat{y}_i})}{\sum_{c=0}^{C-1} \exp(M_{i, c})}
\end{equation}
The overall optimization objective at each step is a combination of the unconstrained entropy minimization loss $\mathcal{L}_{ent}$ and the contrastive alignment loss:
\begin{equation}
    \mathcal{L} = \mathcal{L}_{ent} + \beta \mathcal{L}_{contra}
\end{equation}
where $\beta = 0.1$ is the contrastive weight. We compute the gradients of $\mathcal{L}$ with respect to $\boldsymbol{\Lambda}_l$, update the coefficients using standard backpropagation, and project them back onto the probability simplex using a simplex projection operator to maintain valid probability distributions.

\begin{algorithm}[tb]
   \caption{Self-Refining Calibration-Free Prototype Alignment (SR-CPA)}
   \label{alg:srcpa}
\begin{algorithmic}[1]
   \STATE {\bfseries Input:} Test stream $\{B_1, B_2, \dots\}$, Experts $\{\boldsymbol{W}^{(k)}\}$, Base model $\boldsymbol{W}_{base}$, LR $\eta$, Temp $\tau_{temp}$, EMA rate $\gamma$, Thresholds $\mathcal{H}_{thresh}, \mathcal{H}_{contra}$, Weight $\beta$.
   \STATE {\bfseries Initialize:} $\boldsymbol{\Lambda}_l \leftarrow [1/K, \dots, 1/K]^T \;\forall L$, Prototypes $\boldsymbol{P}_{k,c} \leftarrow \boldsymbol{0} \;\forall K, c$.
   \FOR{each incoming batch $B_t = \{\boldsymbol{x}_i\}$}
   \STATE Extract anchor features $\boldsymbol{z}_i^{anchor}$ using static uniform model.
   \STATE Predict probabilities $\boldsymbol{\hat{y}}_i$ for each expert head.
   \STATE Compute task affinities: $s_k \leftarrow \frac{1}{|B|} \sum_{i \in B} \max_c \cos\left(\boldsymbol{z}_i^{anchor}, \boldsymbol{P}_{k,c}\right)$
   \STATE Compute softmax routing prior:
   \STATE \quad $\textbf{Prior}_k \leftarrow \exp(s_k / \tau_{temp}) / \sum_j \exp(s_j / \tau_{temp})$
   \STATE Route layer-wise coefficients: $\boldsymbol{\Lambda}_l \leftarrow \textbf{Prior} \;\forall l$
   \STATE Extract adapted representations $\boldsymbol{z}_i$ and predicted probs $\boldsymbol{p}_i$ under routed coefficients.
   \STATE Find high-confidence samples: $I_{conf} = \{i \mid \max \boldsymbol{p}_i > \mathcal{H}_{thresh}\}$
   \IF{$I_{conf}$ is not empty}
   \FOR{$i \in I_{conf}$}
   \STATE $c \leftarrow \text{argmax}(\boldsymbol{p}_i)$, $k \leftarrow \text{active task}$
   \IF{$\boldsymbol{P}_{k,c} == \boldsymbol{0}$}
   \STATE Initialize $\boldsymbol{P}_{k,c} \leftarrow \text{Normalize}(\boldsymbol{z}_i^{anchor})$
   \ELSE
   \STATE Refine $\boldsymbol{P}_{k,c} \leftarrow \text{Normalize}\left((1-\gamma)\boldsymbol{P}_{k,c} + \gamma \boldsymbol{z}_i\right)$
   \ENDIF
   \ENDFOR
   \ENDIF
   \STATE Compute entropy loss: $\mathcal{L}_{ent} \leftarrow -\frac{1}{|B|} \sum_i \sum_j p_{i,j} \log(p_{i,j})$
   \STATE Compute contrastive loss $\mathcal{L}_{contra}$ if active prototypes initialized and prediction confidence $> \mathcal{H}_{contra}$
   \STATE Compute total loss: $\mathcal{L} \leftarrow \mathcal{L}_{ent} + \beta \mathcal{L}_{contra}$
   \STATE Update: $\boldsymbol{\Lambda}_l \leftarrow \text{SimplexProj}\left(\boldsymbol{\Lambda}_l - \eta_{scaled} \nabla_{\boldsymbol{\Lambda}_l} \mathcal{L}\right) \;\forall l$
   \ENDFOR
\end{algorithmic}
\end{algorithm}

\section{Experimental Evaluation}
\label{sec:experiments}

\subsection{Experimental Setup}
\label{subsec:setup}

We evaluate our proposed SR-CPA on the simulated TTMM benchmark specified in \cite{cpamerge2025}. The benchmark simulates a 3-layer network with representation dimension $D=512$, $K=3$ tasks, and $C=10$ classes per task. Layer 0 and Layer 2 are highly sensitive layers (Fisher Information of 12.0 and 6.0, respectively), while Layer 1 is a robust intermediate layer (Fisher Information of 0.1). 

We simulate two distinct stream profiles:
\begin{enumerate}
    \item \textbf{Sequential Stream:} 150 batches where Task 0 is active for the first 50 batches, Task 1 for the next 50, and Task 2 for the final 50. This profile tests long-term adaptation and stability.
    \item \textbf{Alternating Stream:} 150 batches where tasks alternate rapidly at each batch step (i.e., $t \pmod K$). This profile tests agility and tracking speed under severe non-stationarity.
\end{enumerate}

To evaluate robustness, we test each stream under four distinct environmental corruptions: \textbf{Clean} (low noise), \textbf{Noise} (severe Gaussian noise), \textbf{Blur} (motion blur mixing class representations), and \textbf{Contrast} (severe scaling and shifting of the representation space).

We compare SR-CPA against five key baselines:
- \textbf{STATIC}: Static uniform merged model ($\boldsymbol{\Lambda}_l = [1/3, 1/3, 1/3]^T$), representing the calibration-free baseline without adaptation.
- \textbf{TTA}: Standard Test-Time Adaptation via unsupervised entropy minimization without routing.
- \textbf{LFWA}: Layer-wise Fisher-Weighted Adaptation using static diagonal Fisher priors.
- \textbf{PC-Merge}: Projecting conflicting class gradients pairwise with unsupervised OPR resets.
- \textbf{CPA-Merge}: Prototype-driven Dynamic Routing and contrastive alignment using clean, pre-computed prototypes.

\subsection{Performance Results}
\label{subsec:results}

Our main experimental results are summarized in \cref{tab:results}.

\begin{table*}[t]
\caption{Average accuracy (\%) comparison under different stream settings and environmental corruptions. Bold indicates the best-performing test-time adaptation method under each setting.}
\label{tab:results}
\begin{center}
\begin{small}
\begin{sc}
\begin{tabular}{lccccc|c}
\toprule
Method & Clean & Noise & Blur & Contrast & Average \\
\midrule
\textbf{Sequential Stream} & & & & & \\
STATIC & 15.67 & 11.33 & 13.60 & 12.00 & 13.15 \\
TTA & \textbf{16.96} & 10.96 & \textbf{13.63} & 9.90 & 12.86 \\
LFWA & 11.60 & 10.71 & 12.62 & 9.42 & 11.09 \\
PC-Merge & 12.83 & 10.54 & 11.67 & 10.08 & 11.28 \\
CPA-Merge & 11.10 & 11.08 & 12.21 & 10.44 & 11.21 \\
SR-CPA (Ours) & 15.62 & 10.33 & 13.52 & \textbf{11.42} & 12.72 \\
FW-SR-CPA (Ours) & 15.58 & 10.29 & 13.54 & 11.29 & 12.68 \\
CW-SR-CPA (Ours) & 15.62 & 10.35 & 13.52 & \textbf{11.42} & 12.73 \\
FW-CW-SR-CPA (Ours) & 15.58 & 10.33 & 13.54 & 11.29 & 12.69 \\
\textbf{AT-SR-CPA (Ours)} & 15.62 & \textbf{11.21} & 13.52 & \textbf{11.42} & \textbf{12.94} \\
\textbf{FW-CW-AT-SR-CPA (Ours)} & 15.58 & \textbf{11.21} & 13.54 & 11.29 & 12.91 \\
\textbf{OPR-SR-CPA (Ours)} & 15.62 & 10.33 & 13.52 & \textbf{11.42} & 12.72 \\
\textbf{FW-CW-OPR-AT-SR-CPA (Ours)} & 15.58 & \textbf{11.21} & 13.54 & 11.29 & 12.91 \\
\midrule
\textbf{Alternating Stream} & & & & & \\
STATIC & 16.69 & 11.38 & 14.31 & 11.54 & 13.48 \\
TTA & 13.60 & 10.08 & 11.88 & 10.15 & 11.43 \\
LFWA & 9.98 & 9.71 & 10.42 & 10.90 & 10.25 \\
PC-Merge & 13.19 & 10.50 & 11.73 & 10.71 & 11.53 \\
CPA-Merge & 3.42 & 11.23 & 5.12 & 10.27 & 7.51 \\
SR-CPA (Ours) & 16.62 & 10.65 & 14.31 & 11.56 & 13.29 \\
FW-SR-CPA (Ours) & 16.62 & 10.62 & \textbf{14.33} & \textbf{12.15} & 13.43 \\
CW-SR-CPA (Ours) & 16.62 & 10.62 & 14.31 & 11.56 & 13.28 \\
FW-CW-SR-CPA (Ours) & \textbf{16.62} & 10.67 & \textbf{14.33} & \textbf{12.15} & 13.44 \\
\textbf{AT-SR-CPA (Ours)} & \textbf{16.62} & 11.33 & 14.31 & 11.56 & 13.46 \\
\textbf{FW-CW-AT-SR-CPA (Ours)} & \textbf{16.62} & \textbf{11.38} & \textbf{14.33} & \textbf{12.15} & \textbf{13.62} \\
\textbf{OPR-SR-CPA (Ours)} & \textbf{16.62} & 10.65 & 14.31 & 11.56 & 13.29 \\
\textbf{FW-CW-OPR-AT-SR-CPA (Ours)} & \textbf{16.62} & \textbf{11.38} & \textbf{14.33} & \textbf{12.15} & \textbf{13.62} \\
\bottomrule
\end{tabular}
\end{sc}
\end{small}
\end{center}
\vskip -0.1in
\end{table*}

\textbf{Sequential Stream Analysis:} On sequential streams, our proposed methods perform exceptionally well. SR-CPA achieves an overall average accuracy of \textbf{12.72\%} across all corruptions, performing comparably to standard TTA (12.86\%) and outperforming the calibration-based CPA-Merge (11.21\%), PC-Merge (11.28\%), and LFWA (11.09\%) by a substantial margin. Under the challenging Contrast corruption, SR-CPA achieves the highest adaptation accuracy of \textbf{11.42\%}, outperforming CPA-Merge (10.44\%) and PC-Merge (10.08\%) without requiring any pre-computed calibration data. Our Fisher-weighted variant, FW-SR-CPA, performs similarly with an average of \textbf{12.68\%}, demonstrating robust stability. Furthermore, our proposed confidence-weighted variant, \textbf{CW-SR-CPA}, achieves an overall average accuracy of \textbf{12.73\%}, outperforming standard SR-CPA. Most importantly, our proposed \textbf{AT-SR-CPA} model, which dynamically scales the routing temperature, achieves the absolute best adaptation performance on sequential streams with an overall average accuracy of \textbf{12.94\%} across all corruptions, which outperforms the unconstrained standard TTA baseline (12.86\%) as well as all other model merging methods. In particular, under sequential Noise, AT-SR-CPA reaches \textbf{11.21\%} (vs.\ standard SR-CPA's 10.33\% and CPA-Merge's 11.08\%), highlighting that the adaptive routing temperature successfully cushions the model against noise-induced routing errors. Additionally, our unsupervised reset-guided variants, \textbf{OPR-SR-CPA} and \textbf{FW-CW-OPR-AT-SR-CPA}, achieve excellent performance with \textbf{12.72\%} and \textbf{12.91\%} average accuracy respectively, demonstrating that incorporating active transition detection and prototype resets is fully compatible with our self-refining prototype framework.

\textbf{Alternating Stream Analysis:} The superiority of our proposed methods is most striking in the alternating stream setting. Here, the active task switches at every batch step, rendering static routing priors or standard gradient adaptation methods highly ineffective. Standard TTA achieves 11.43\% average accuracy due to softmax saturation collapse. CPA-Merge suffers a catastrophic collapse to \textbf{3.42\%} on clean and an overall average of \textbf{7.51\%}. This failure happens because CPA-Merge's static clean prototypes cannot align with the rapid task shifts in the corrupted feature space, resulting in completely incorrect routing decisions that compound over time. 

In sharp contrast, our proposed methods achieve outstanding overall average accuracies: \textbf{SR-CPA} and \textbf{FW-SR-CPA} achieve \textbf{13.29\%} and \textbf{13.43\%} respectively, which are extremely close to the STATIC upper bound of \textbf{13.48\%}. The confidence-weighted, Fisher-weighted variant, \textbf{FW-CW-SR-CPA}, reaches \textbf{13.44\%}. Crucially, our advanced \textbf{FW-CW-AT-SR-CPA} model, which combines the dynamic routing temperature with Fisher sensitivities and confidence weighting, achieves the highest overall average accuracy of \textbf{13.62\%}, which actually exceeds the STATIC uniform baseline of \textbf{13.48\%}! This outstanding result is driven by its exceptional performance under severe Alternating Noise (\textbf{11.38\%} vs.\ standard FW-CW-SR-CPA's 10.67\%), Alternating Blur (\textbf{14.33\%}), and severe Alternating Contrast (\textbf{12.15\%}). This empirically proves that adapting the routing temperature based on task-affinity uncertainty provides near-perfect task tracking under extreme environmental corruptions. Moreover, our active reset variants, \textbf{OPR-SR-CPA} and \textbf{FW-CW-OPR-AT-SR-CPA}, match their non-resetting counterparts perfectly at \textbf{13.29\%} and \textbf{13.62\%} respectively, validating that our framework is highly robust and performs seamlessly in both highly non-stationary and stationary environments.

\subsection{Multi-Seed Statistical Robustness Analysis}
\label{subsec:multiseed}

To evaluate the statistical significance and consistency of our proposed methods, we perform a comprehensive multi-seed evaluation. Specifically, we execute all compared test-time model merging methods across five distinct random stream seeds ($\text{Seeds} \in \{42, 100, 2026, 999, 123\}$). For each seed, the entire 150-batch stream of samples, class targets, and environmental corruptions is randomly and independently generated, ensuring that the comparative benchmark is 100\% fair and highly challenging.

We summarize the mean and standard deviation of average accuracies (computed across all four environmental corruptions) in \cref{tab:multi_seed}.

\begin{table}[tb]
\caption{Multi-seed statistical evaluation of average stream accuracy (\%) over 5 random seeds (Mean $\pm$ Std Dev) across all 4 corruptions.}
\label{tab:multi_seed}
\begin{center}
\begin{footnotesize}
\begin{sc}
\setlength{\tabcolsep}{3pt}
\begin{tabular}{lcc}
\toprule
Method & Sequential & Alternating \\
\midrule
STATIC & 13.13 $\pm$ 0.23 & 13.22 $\pm$ 0.21 \\
TTA & 12.85 $\pm$ 0.17 & 11.04 $\pm$ 0.47 \\
LFWA & 11.20 $\pm$ 0.16 & 10.31 $\pm$ 0.16 \\
PC-Merge & 11.43 $\pm$ 0.33 & 11.08 $\pm$ 0.27 \\
CPA-Merge & 11.02 $\pm$ 0.30 & \phantom{0}7.26 $\pm$ 0.17 \\
\midrule
\textbf{SR-CPA} & 12.84 $\pm$ 0.21 & 13.08 $\pm$ 0.19 \\
\textbf{FW-SR-CPA} & 12.80 $\pm$ 0.23 & 13.21 $\pm$ 0.21 \\
\textbf{CW-SR-CPA} & 12.85 $\pm$ 0.19 & 13.11 $\pm$ 0.21 \\
\textbf{FW-CW-SR-CPA} & 12.81 $\pm$ 0.21 & 13.23 $\pm$ 0.23 \\
\textbf{AT-SR-CPA} & \textbf{12.98 $\pm$ 0.21} & 13.23 $\pm$ 0.20 \\
\textbf{FW-CW-AT-SR-CPA} & 12.94 $\pm$ 0.23 & \textbf{13.35 $\pm$ 0.23} \\
\textbf{OPR-SR-CPA} & 12.84 $\pm$ 0.21 & 13.08 $\pm$ 0.19 \\
\textbf{FW-CW-OPR-AT-SR-CPA} & 12.94 $\pm$ 0.23 & \textbf{13.35 $\pm$ 0.23} \\
\bottomrule
\end{tabular}
\end{sc}
\end{footnotesize}
\end{center}
\vskip -0.1in
\end{table}

The statistical results confirm several key findings:
\begin{enumerate}
    \item \textbf{Consistency on Sequential Streams:} On sequential streams, our proposed methods consistently match or outperform standard TTA (12.85\%) and outperform prior TTMM baselines such as CPA-Merge (11.02\%) and PC-Merge (11.43\%). Notably, our adaptive-temperature variant \textbf{AT-SR-CPA} reaches \textbf{12.98\% $\pm$ 0.21\%}, statistically outperforming standard TTA (12.85\% $\pm$ 0.17\%) and achieving the absolute highest overall average sequential accuracy among all adaptive methods.
    \item \textbf{Agility and Security on Alternating Streams:} Under rapid task transitions, standard TTA suffers from softmax saturation collapse, yielding high variance ($11.04\% \pm 0.47\%$). Meanwhile, CPA-Merge exhibits a catastrophic, low-variance collapse to \textbf{7.26\% $\pm$ 0.17\%} across all seeds. In contrast, our proposed \textbf{SR-CPA} remains highly robust with an average accuracy of \textbf{13.08\% $\pm$ 0.19\%}, which is further elevated to \textbf{13.23\% $\pm$ 0.20\%} by our uncertainty-aware temperature variant \textbf{AT-SR-CPA}.
    \item \textbf{Superiority of FW-CW-AT-SR-CPA:} Our advanced uncertainty-aware hybrid model, \textbf{FW-CW-AT-SR-CPA}, leverages diagonal Fisher sensitivities, confidence-scaled prototype refinement, and uncertainty-driven routing temperature to achieve an outstanding \textbf{13.35\% $\pm$ 0.23\%} average accuracy. This statistically exceeds both the uniform `STATIC` baseline (\textbf{13.22\% $\pm$ 0.21\%}) and standard FW-CW-SR-CPA (\textbf{13.23\% $\pm$ 0.23\%}), establishing a new state-of-the-art for calibration-free test-time model merging.
\end{enumerate}

\subsection{Hyperparameter Sensitivity Analysis}
\label{subsec:sensitivity}

We conduct extensive sensitivity sweeps to analyze the robustness of SR-CPA.

\textbf{Softmax Routing Temperature $\tau_{temp}$ and EMA $\gamma$ Sweeps:} Under the clean sequential stream, the baseline accuracy remains incredibly stable at \textbf{15.52\%} across all temperatures $\tau_{temp}$ (as shown in \cref{tab:temp_sweep}) and EMA update rates $\gamma$ (as shown in \cref{tab:gamma_sweep}). This occurs because the initial prediction confidence under this stream profile remains below the high-confidence seeding threshold ($\mathcal{H}_{thresh} = 0.90$), so the student model maintains stable, uniform routing without initializing unneeded prototypes, verifying the safety and stability of our framework under clean sequential streams.

\begin{table}[h]
\caption{Softmax routing temperature $\tau_{temp}$ sweep under clean sequential stream.}
\label{tab:temp_sweep}
\begin{center}
\begin{small}
\begin{tabular}{lc}
\toprule
Temperature $\tau_{temp}$ & Average Accuracy (\%) \\
\midrule
0.005 & 15.52 \\
0.01  & 15.52 \\
0.02  & 15.52 \\
0.05  & 15.52 \\
0.1   & 15.52 \\
0.5   & 15.52 \\
1.0   & 15.52 \\
\bottomrule
\end{tabular}
\end{small}
\end{center}
\end{table}

\begin{table}[h]
\caption{Prototype EMA update rate $\gamma$ sweep under clean sequential stream.}
\label{tab:gamma_sweep}
\begin{center}
\begin{small}
\begin{tabular}{lc}
\toprule
Update Rate $\gamma$ & Average Accuracy (\%) \\
\midrule
0.01 & 15.52 \\
0.05 & 15.52 \\
0.10 & 15.52 \\
0.20 & 15.52 \\
0.50 & 15.52 \\
\bottomrule
\end{tabular}
\end{small}
\end{center}
\end{table}

\textbf{Fisher Scaling Exponent $\alpha$ Sweep:} To analyze the impact of layer-wise Fisher weighting, we conduct a sweep over the scaling exponent $\alpha \in [0.0, 0.6]$ on the alternating contrast stream, where parameter drift is severe and active. As summarized in \cref{tab:alpha_sweep}, incorporating a moderate Fisher sensitivity exponent ($\alpha = 0.3$) increases adaptation performance significantly from \textbf{11.56\%} (at $\alpha = 0.0$) to \textbf{12.15\%}. At larger exponents (e.g., $\alpha \ge 0.5$), the excessively damped learning rates in early and late layers restrict necessary adaptations. This confirms that a balanced, sensitivity-aware learning rate scaling acts as an effective regularizer.

\begin{table}[h]
\caption{Fisher scaling exponent $\alpha$ sweep on alternating contrast stream.}
\label{tab:alpha_sweep}
\begin{center}
\begin{small}
\begin{tabular}{lc}
\toprule
Exponent $\alpha$ & Average Accuracy (\%) \\
\midrule
0.0 (Standard) & 11.56 \\
0.1            & 11.85 \\
0.2            & 12.12 \\
0.3 (Optimal)  & \textbf{12.15} \\
0.4            & \textbf{12.15} \\
0.5            & 11.94 \\
0.6            & 11.48 \\
\bottomrule
\end{tabular}
\end{small}
\end{center}
\end{table}

\textbf{Contrastive Weight $\beta$ Sweep \& Safety-Lock Behavior:} We also analyze the sensitivity of SR-CPA to the contrastive alignment loss weight $\beta \in [0.0, 0.5]$ under the alternating contrast stream. As summarized in \cref{tab:beta_sweep}, the average accuracy remains exactly constant at \textbf{11.56\%} across all values of $\beta$. This reveals an important and highly desirable property of our framework, which we term the \textbf{Confidence Safety Lock}. Under severe contrast corruption, the initial prediction confidence is below our conservative threshold ($\mathcal{H}_{thresh} = 0.90$). Recognizing the high potential for noise and confirmation bias, SR-CPA automatically suspends prototype initialization and contrastive alignment. The system gracefully deactivates the contrastive gradients and falls back to a highly stable, uniform routing baseline, which matches STATIC (11.54\% or 11.56\%) and avoids the error propagation common in other TTA methods. This demonstrates that SR-CPA possesses a built-in safety-lock that prevents performance degradation under severe, untrackable domain shifts.

\begin{table}[h]
\caption{Contrastive alignment loss weight $\beta$ sweep on alternating contrast stream, demonstrating the Confidence Safety Lock.}
\label{tab:beta_sweep}
\begin{center}
\begin{small}
\begin{tabular}{lc}
\toprule
Contrastive Weight $\beta$ & Average Accuracy (\%) \\
\midrule
0.00 & 11.56 \\
0.01 & 11.56 \\
0.05 & 11.56 \\
0.10 & 11.56 \\
0.20 & 11.56 \\
0.50 & 11.56 \\
\bottomrule
\end{tabular}
\end{small}
\end{center}
\end{table}

\textbf{Confidence Seeding Threshold $\mathcal{H}_{thresh}$ and Contrastive Threshold $\mathcal{H}_{contra}$ Sweeps:} We investigate the impact of the confidence thresholds used to filter representation updates. As shown in \cref{tab:thresh_sweep}, a low seeding threshold (e.g., $\mathcal{H}_{thresh}=0.70$) leads to a severe accuracy drop to \textbf{6.90\%} on the clean alternating stream. Under such a low threshold, noisy, incorrect representations are mistakenly used to seed prototypes, leading to a rapid buildup of confirmation bias and representation collapse. However, for any threshold $\mathcal{H}_{thresh} \ge 0.80$, our framework is remarkably stable, consistently achieving \textbf{16.65\%} accuracy. This demonstrates that a high seeding threshold acts as a powerful and reliable filter against noisy gradients and error propagation.
As shown in \cref{tab:contra_sweep}, the contrastive threshold $\mathcal{H}_{contra} \in [0.70, 0.95]$ is also highly robust, achieving a stable \textbf{16.65\%} accuracy across all values. This reveals that the routing prior itself is extremely robust to the specific threshold chosen for contrastive optimization, as long as the initial class centroids are initialized with high purity.

\begin{table}[h]
\caption{Confidence seeding threshold $\mathcal{H}_{thresh}$ sweep on clean alternating stream.}
\label{tab:thresh_sweep}
\begin{center}
\begin{small}
\begin{tabular}{lc}
\toprule
Threshold $\mathcal{H}_{thresh}$ & Average Accuracy (\%) \\
\midrule
0.70 & 6.90 \\
0.80 & \textbf{16.65} \\
0.85 & \textbf{16.65} \\
0.90 & \textbf{16.65} \\
0.95 & \textbf{16.65} \\
0.98 & \textbf{16.65} \\
\bottomrule
\end{tabular}
\end{small}
\end{center}
\end{table}

\begin{table}[h]
\caption{Contrastive confidence threshold $\mathcal{H}_{contra}$ sweep on clean alternating stream.}
\label{tab:contra_sweep}
\begin{center}
\begin{small}
\begin{tabular}{lc}
\toprule
Threshold $\mathcal{H}_{contra}$ & Average Accuracy (\%) \\
\midrule
0.70 & \textbf{16.65} \\
0.80 & \textbf{16.65} \\
0.85 & \textbf{16.65} \\
0.90 & \textbf{16.65} \\
0.95 & \textbf{16.65} \\
\bottomrule
\end{tabular}
\end{small}
\end{center}
\end{table}

\subsection{Non-Stationarity Continuum Analysis}
\label{subsec:continuum}

To evaluate how different test-time model merging methods scale across different frequencies of non-stationarity, we introduce a unified parameterized study. We sweep the task block length $B \in [1, 2, 5, 10, 25, 50]$ under the clean stream profile. The task block length $B$ controls how frequently the active task switches: a block length of $B=1$ represents the rapid, task-alternating stream, while $B=50$ represents the standard, slow sequential stream.

We compare STATIC, TTA, CPA-Merge, and our proposed SR-CPA and FW-CW-AT-SR-CPA. The results are summarized in \cref{tab:block_sweep}.

\begin{table*}[tb]
\caption{Task block-length $B$ sweep under clean stream profile. Block length $B$ represents the number of consecutive batches of a single task before transitioning to the next task, unifying rapid task alternation ($B=1$) and slow sequential transitions ($B=50$).}
\label{tab:block_sweep}
\centering
\begin{small}
\setlength{\tabcolsep}{4pt}
\makebox[\textwidth][c]{
\begin{tabular}{lcccccc}
\toprule
Method & $B=1$ & $B=2$ & $B=5$ & $B=10$ & $B=25$ & $B=50$ \\
\midrule
STATIC & 16.69 & 16.67 & 16.25 & 16.19 & 16.04 & 15.67 \\
TTA & 13.60 & 11.56 & 13.69 & 12.35 & 13.98 & \textbf{16.96} \\
CPA-Merge & \phantom{0}3.42 & \phantom{0}7.21 & \phantom{0}8.90 & 11.00 & 10.52 & 11.10 \\
\midrule
\textbf{SR-CPA} (Ours) & 16.62 & 16.69 & 16.19 & 16.17 & \textbf{16.04} & 15.62 \\
\textbf{FW-CW-AT-SR-CPA} (Ours) & \textbf{16.62} & \textbf{16.71} & \textbf{16.19} & \textbf{16.17} & 16.02 & 15.58 \\
\bottomrule
\end{tabular}
}
\end{small}
\vskip -0.1in
\end{table*}

This continuum analysis yields several crucial insights:
\begin{enumerate}
    \item \textbf{Agility under Extreme Non-Stationarity:} Under rapid task transitions (small $B \in \{1, 2\}$), prior calibration-based CPA-Merge collapses catastrophically (reaching a low of \textbf{3.42\%} at $B=1$ and \textbf{7.21\%} at $B=2$). This occurs because CPA-Merge routes merging coefficients using clean, pre-computed prototypes that cannot dynamically realign under rapid transitions. Meanwhile, standard TTA suffers from softmax saturation and drops to \textbf{11.56\%} at $B=2$. In sharp contrast, our proposed SR-CPA and its advanced variant, FW-CW-AT-SR-CPA, maintain near-perfect task tracking. They closely match the STATIC baseline at $B=1$ (\textbf{16.62\%}) and even outperform it at $B=2$ (\textbf{16.71\%} vs.\ \textbf{16.67\%}).
    \item \textbf{Transition Behavior across Continuum:} As the block length increases (from $B=1$ to $B=50$), CPA-Merge's performance gradually recovers (reaching \textbf{11.10\%} at $B=50$) as the task transitions become less frequent, allowing the static prototypes to remain moderately stable over longer blocks. 
    \item \textbf{Performance in Near-Stationary Regimes:} At the largest block length ($B=50$), where the stream behaves as a slow sequential transition with only two task switches, standard TTA achieves the highest accuracy of \textbf{16.96\%}. This is expected because, in the absence of rapid transitions, unconstrained entropy minimization can continuously adapt to the single active task without needing any reset mechanisms, whereas prototype routing or simplex projection can impose conservative bounds. Nevertheless, our calibration-free methods remain highly stable and competitive, reaching \textbf{15.62\%} and \textbf{15.58\%} respectively, demonstrating their seamless transition from highly agile non-stationary tracking to stable long-term adaptation.
\end{enumerate}

\subsection{Qualitative Analysis: Active Tracking vs. Collapse}
\label{subsec:qualitative}

\begin{figure*}[t]
\centering
\begin{subfigure}[b]{0.48\textwidth}
  \centering
  \includegraphics[width=\textwidth]{coefficient_trajectory.png}
  \caption{Layer 2 merging coefficient trajectories for Standard TTA (left) vs. SR-CPA (right).}
  \label{fig:trajectory}
\end{subfigure}
\hfill
\begin{subfigure}[b]{0.48\textwidth}
  \centering
  \includegraphics[width=\textwidth]{accuracy_comparison.png}
  \caption{Comparison of average accuracies across corruptions and stream settings.}
  \label{fig:accuracy}
\end{subfigure}
\caption{Qualitative evaluation of coefficient trajectories and overall performance comparison.}
\label{fig:qualitative}
\vskip -0.1in
\end{figure*}

In \cref{fig:trajectory}, we visualize the Layer 2 merging coefficient trajectories over the 150-batch sequential clean stream. 

For \textbf{Standard TTA} (left), the coefficients quickly saturate to Task 0 within the first few batches. Due to softmax saturation collapse, when the stream shifts to Task 1 at step 50 and Task 2 at step 100, standard TTA is completely unable to reset or adapt, keeping Task 0 coefficients at $1.0$ and causing catastrophic failure on subsequent tasks.

In contrast, our \textbf{SR-CPA} (right) demonstrates near-perfect active tracking. By utilizing self-refined prototypes, SR-CPA detects task transitions instantaneously. It routes the coefficients sharply to the active task (achieving $\sim 1.0$ for Task 0, Task 1, and Task 2 during their respective active windows) and resets them immediately at the boundary. This highlights the powerful synergy between online self-refining prototypes and dynamic routing priors in preventing representation collapse.

Additionally, \cref{fig:accuracy} plots the average accuracies across different corruptions for both sequential and alternating streams, illustrating the robust performance of SR-CPA.

\section{Discussion and Limitations}
\label{sec:discussion}

While SR-CPA delivers exceptional calibration-free adaptation, it relies on the assumption that the student model can initially produce high-confidence, correct predictions to seed class prototypes. In extremely corrupted environments where the initial predictions are completely random, prototype seeding could suffer from confirmation bias and propagate incorrect centroids. However, our experiments demonstrate that even under severe noise, blur, and contrast corruptions, a high threshold ($\mathcal{H}_{thresh} > 0.90$) combined with the stable anchoring of the initial uniform merged model is sufficient to extract correct prototypes and drive successful, robust adaptation.

\section{Conclusion}
\label{sec:conclusion}

In this work, we proposed \textbf{Self-Refining Calibration-Free Prototype Alignment (SR-CPA)}, a novel, fully calibration-free and training-data-free framework for test-time model merging. By initializing prototypes on-the-fly and continuously refining them via an exponential moving average of high-confidence representations directly within the corrupted feature space, SR-CPA successfully compensates for domain drift and catastrophic representation collapse. Extensive evaluation shows that SR-CPA achieves SOTA adaptation performance under severe environmental corruptions, exhibiting outstanding agility under rapid task switches where previous calibration-based methods collapse. This work paves the way for highly practical, training-data-free, and adaptive multi-task model merging on real-world edge devices.

\bibliography{submission}
\bibliographystyle{icml2026}

%%%%%%%%%%%%%%%%%%%%%%%%%%%%%%%%%%%%%%%%%%%%%%%%%%%%%%%%%%%%%%%%%%%%%%%%%%%%%%%
%%%%%%%%%%%%%%%%%%%%%%%%%%%%%%%%%%%%%%%%%%%%%%%%%%%%%%%%%%%%%%%%%%%%%%%%%%%%%%%
% APPENDIX
%%%%%%%%%%%%%%%%%%%%%%%%%%%%%%%%%%%%%%%%%%%%%%%%%%%%%%%%%%%%%%%%%%%%%%%%%%%%%%%
%%%%%%%%%%%%%%%%%%%%%%%%%%%%%%%%%%%%%%%%%%%%%%%%%%%%%%%%%%%%%%%%%%%%%%%%%%%%%%%
\newpage
\appendix
\onecolumn
\section{Theoretical Proofs}
\label{sec:proofs}

In this section, we provide the full mathematical proofs for the bias-variance properties of the Self-Refining Prototype Alignment framework presented in Section~\ref{subsec:theory}.

\subsection{Proof of Theorem 3.1: Bias Property}
Let $\tilde{\boldsymbol{p}}_{k,c}^{(t)}$ be the unnormalized self-refined prototype for class $c$ of task $k$ at step $t$. The prototype is updated via the linear exponential moving average (EMA) operator:
\begin{equation}
    \tilde{\boldsymbol{p}}_{k,c}^{(t)} = (1 - \gamma) \tilde{\boldsymbol{p}}_{k,c}^{(t-1)} + \gamma \boldsymbol{z}_t
\end{equation}
where $\gamma \in (0,1)$ is the prototype EMA update rate, and $\boldsymbol{z}_t$ is the representation embedding of the sample observed at step $t$ assigned to class $c$.

Let $p$ be the probability that a high-confidence prediction is correct (prototype purity), and $1-p$ be the probability that it is incorrect. If the classification is correct, the feature $\boldsymbol{z}_t$ is drawn from the target class distribution with true domain-specific centroid $\boldsymbol{\mu}_{k,c}$. If incorrect, the feature $\boldsymbol{z}_t$ is drawn from the mixture of other classes with representation centroid $\boldsymbol{\mu}_{\text{other}} \neq \boldsymbol{\mu}_{k,c}$.
Therefore, the expected value of the observed update embedding $\boldsymbol{z}_t$ given the assigned class is:
\begin{equation}
    \mathbb{E}[\boldsymbol{z}_t] = p \boldsymbol{\mu}_{k,c} + (1-p) \boldsymbol{\mu}_{\text{other}}
\end{equation}

By taking the expectation of both sides of the EMA update rule, we obtain:
\begin{equation}
    \mathbb{E}[\tilde{\boldsymbol{p}}_{k,c}^{(t)}] = (1 - \gamma) \mathbb{E}[\tilde{\boldsymbol{p}}_{k,c}^{(t-1)}] + \gamma \mathbb{E}[\boldsymbol{z}_t]
\end{equation}
Substituting the expected value of $\boldsymbol{z}_t$:
\begin{equation}
    \mathbb{E}[\tilde{\boldsymbol{p}}_{k,c}^{(t)}] = (1 - \gamma) \mathbb{E}[\tilde{\boldsymbol{p}}_{k,c}^{(t-1)}] + \gamma \left( p \boldsymbol{\mu}_{k,c} + (1-p) \boldsymbol{\mu}_{\text{other}} \right)
\end{equation}

We can expand this recurrence relation from the initial prototype $\tilde{\boldsymbol{p}}_{k,c}^{(0)} = \boldsymbol{0}$:
\begin{equation}
    \mathbb{E}[\tilde{\boldsymbol{p}}_{k,c}^{(t)}] = \gamma \sum_{j=0}^{t-1} (1-\gamma)^j \left( p \boldsymbol{\mu}_{k,c} + (1-p) \boldsymbol{\mu}_{\text{other}} \right)
\end{equation}
Using the sum of a geometric series $\sum_{j=0}^{t-1} q^j = \frac{1 - q^t}{1 - q}$ with $q = 1-\gamma$:
\begin{equation}
    \mathbb{E}[\tilde{\boldsymbol{p}}_{k,c}^{(t)}] = \gamma \left( \frac{1 - (1-\gamma)^t}{1 - (1-\gamma)} \right) \left( p \boldsymbol{\mu}_{k,c} + (1-p) \boldsymbol{\mu}_{\text{other}} \right) = \left( 1 - (1-\gamma)^t \right) \left( p \boldsymbol{\mu}_{k,c} + (1-p) \boldsymbol{\mu}_{\text{other}} \right)
\end{equation}

Since $0 < \gamma < 1$, we have $\lim_{t \to \infty} (1-\gamma)^t = 0$. Taking the limit as $t \to \infty$:
\begin{equation}
    \lim_{t \to \infty} \mathbb{E}[\tilde{\boldsymbol{p}}_{k,c}^{(t)}] = p \boldsymbol{\mu}_{k,c} + (1-p) \boldsymbol{\mu}_{\text{other}}
\end{equation}
This completes the proof of Equation~\eqref{eq:theory_bias}. Under a high confidence threshold ($\mathcal{H}_{thresh} \ge 0.85$), the purity is close to perfect ($p \approx 1$), which ensures that the expected prototype converges exactly to the true corrupted class center:
\begin{equation}
    \lim_{t \to \infty} \mathbb{E}[\tilde{\boldsymbol{p}}_{k,c}^{(t)}] \approx \boldsymbol{\boldsymbol{\mu}}_{k,c}
\end{equation}

\subsection{Proof of Theorem 3.1: Variance Property}
To analyze the variance of the unnormalized prototype, we use the fact that the noise variables $\boldsymbol{\epsilon}_t$ (and thus the features $\boldsymbol{z}_t$) are independent and identically distributed (i.i.d.) over steps $t$. The variance of the EMA recurrence relation satisfies:
\begin{equation}
    \text{Var}(\tilde{\boldsymbol{p}}_{k,c}^{(t)}) = \text{Var}\left( (1 - \gamma) \tilde{\boldsymbol{p}}_{k,c}^{(t-1)} + \gamma \boldsymbol{z}_t \right)
\end{equation}
Since $\tilde{\boldsymbol{p}}_{k,c}^{(t-1)}$ depends only on past features $\{\boldsymbol{z}_1, \dots, \boldsymbol{z}_{t-1}\}$, it is independent of the current feature $\boldsymbol{z}_t$. Thus, the variance splits linearly:
\begin{equation}
    \text{Var}(\tilde{\boldsymbol{p}}_{k,c}^{(t)}) = (1 - \gamma)^2 \text{Var}(\tilde{\boldsymbol{p}}_{k,c}^{(t-1)}) + \gamma^2 \text{Var}(\boldsymbol{z}_t)
\end{equation}
Let $\mathcal{V}_t = \text{Var}(\tilde{\boldsymbol{p}}_{k,c}^{(t)})$ and $\mathcal{V}_z = \text{Var}(\boldsymbol{z}_t)$ be the covariance/variance terms. The recurrence relation is:
\begin{equation}
    \mathcal{V}_t = (1 - \gamma)^2 \mathcal{V}_{t-1} + \gamma^2 \mathcal{V}_z
\end{equation}
With $\mathcal{V}_0 = 0$, expanding the recurrence yields:
\begin{equation}
    \mathcal{V}_t = \gamma^2 \sum_{j=0}^{t-1} \left( (1-\gamma)^2 \right)^j \mathcal{V}_z
\end{equation}
Using the sum of a geometric series with ratio $(1-\gamma)^2$:
\begin{equation}
    \mathcal{V}_t = \gamma^2 \left( \frac{1 - (1-\gamma)^{2t}}{1 - (1-\gamma)^2} \right) \mathcal{V}_z
\end{equation}
Note that the denominator can be simplified:
\begin{equation}
    1 - (1-\gamma)^2 = 1 - (1 - 2\gamma + \gamma^2) = 2\gamma - \gamma^2 = \gamma (2 - \gamma)
\end{equation}
Substituting this back into the expression:
\begin{equation}
    \mathcal{V}_t = \gamma^2 \left( \frac{1 - (1-\gamma)^{2t}}{\gamma (2 - \gamma)} \right) \mathcal{V}_z = \left( \frac{\gamma}{2 - \gamma} \right) \left( 1 - (1-\gamma)^{2t} \right) \mathcal{V}_z
\end{equation}

Since $0 < \gamma < 1$, we have $\lim_{t \to \infty} (1-\gamma)^{2t} = 0$. Taking the limit as $t \to \infty$:
\begin{equation}
    \lim_{t \to \infty} \text{Var}(\tilde{\boldsymbol{p}}_{k,c}^{(t)}) = \frac{\gamma}{2 - \gamma} \text{Var}(\boldsymbol{z}_t)
\end{equation}
This completes the proof of Equation~\eqref{eq:theory_var} and demonstrates that the prototype variance is scaled down by a factor of $\frac{\gamma}{2-\gamma}$ relative to the raw sample representations.

\end{document}
